% SY52210A
% Sauerstoffberieselung
% 14.0
% 2013-11-30

.. _m-sauerstoffberieselung:


\label{proc:SY52210A}
Bei jedem Patienten, bei dem eine
lebensbedrohliche Störung einer vitalen Funktion eingetreten
ist oder einzutreten droht (\enquote{Notfallpatient}), soll,
sofern keine Kontraindiktaionen vorliegen, soviel Sauerstoff
verabreicht werden, sodass die Sauerstoffsättigung
(SpO₂) im Bereich von \EE{\VonBis{94}{98}%} erreicht wird.

\begin{enumerate}
  -   Kontraindikationen und Gegenanzeigen prüfen:
  \begin{itemize}
    -   COPD (\FullPrettyRef{topic:copd})
    -   Hyperventilationssyndrom, Hyperventilationstetanie
    (\FullPrettyRef{topic:hyperventilationssyndrom})
  \end{itemize}
  -   Situationsgerechte Dosierung je nach
  zugrundeliegender Erkrankung. Grundsätzlich soll ein
  SpO₂ von
\EE{\VonBis{94}{98}%} erreicht werden. Steht keine Pulsoxymetrie zur Verfügung, ist als
Richtwert von einer Dosis von \EE{8 L / min} auszugehen,
welche dem klinischen Zustand des Patienten angepasst werden
muss.
  -   Auswahl des Hilfsmittels:
  \begin{itemize}
    -   \E{Sauerstoff\-brille}
    O₂-Flow bis~5 L / min,
    -   \E{Sauerstoff\-maske}
    O₂-Flow ab 5 L / min,
    -   \E{Sauerstoff\-maske mit
    Reservoir} O₂-Flow ab 5 L / min,
  \end{itemize}
  -   Patientenaufklärung
  -   Sauerstoffgerät einschalten und Durchflußrate
  einstellen
  -   Bei Verwendung einer Sauerstoff\-maske mit
  Reservoir: Reservoir füllen
  -   Sauerstoffbrille/-maske positionieren 
\end{enumerate} 

:cite:`AsboeLehrmeinungRd2011-09`\ 